\clearpage\thispagestyle{empty}
\pdfbookmark[0]{Quellen zum Brief über Grahams Planeten vom 4. Juni 1848}{graham641-sources}
\begin{center}\textbf{Quellen zum Brief über Grahams Planeten vom 4. Juni 1848}\end{center}
\noindent Diese KI-Übertragung ins heutige Deutsch folgt dem Brief auf S.~477 der \href{https://gdz.sub.uni-goettingen.de/id/PPN236007467}{\textit{Werke}, Band~VI} (Göttingen, 1874). Carl Friedrich Gauß schrieb am 4. Juni 1848 an Heinrich Christian Schumacher. Er fügt eine Beobachtung vom 24. Mai hinzu und gibt die von Goldschmidt berechneten Elemente an.
\par\smallskip\noindent Der Brief erschien zuerst in den \textit{Astronomischen Nachrichten}, Band~XXVII, Nr.~641, \href{https://api.digitale-sammlungen.de/iiif/image/v2/bsb10049452_00141/full/1800,/0/default.jpg}{Sp.~265--266}. Die \href{https://api.digitale-sammlungen.de/iiif/image/v2/bsb10049452_00144/full/1800,/0/default.jpg}{Schlusszeile auf Sp.~272} nennt den 29. Juni 1848. Die Zeitschrift druckt unmittelbar nach dem Brief einen Vergleich der Meridianbeobachtungen mit den Elementen in 27 Zeilen ab. Der Text auf S.~477 endet mit dem Satz über die Abnahme des Instruments. KI-Übertragung und KI-Satz.
\clearpage
